% GENERATED FILE. Edit canonical lessons, then run npm run generate:books.
% canonical-source-hash: fnv1a64:9535622dd384c88f
% canonical-lessons: BN-C444-obhiman, BN-C444-onujog, BN-C444-mukh-bhar-kora, BN-C444-khoma-kora, BN-C444-monomalinyo

\chapter{Sulking and Making Up}
\label{ch:bn-sulking-and-making-up}
\hlchaptermodality{444}

\begin{chapteropening}
\textbf{By the end of this chapter:} \emph{I can talk about sulking, misunderstandings and making up.}

Complete the last lesson of chapter 444: I can talk about sulking, misunderstandings and making up.
\end{chapteropening}

\section[\texorpdfstring{\bn{অভিমান} (ôbhimān)}{অভিমান (ôbhimān)}]{\bn{অভিমান} (ôbhimān) — hurt pride; a loving sulk}

\label{lesson:BN-C444-obhiman}



\begin{quote}
Before the new one: say the Bengali for to seethe, then the Bengali for to grumble.
\end{quote}

\subsection*{You'll want to know: \bn{অভিমান}}
\textbf{\bn{অভিমান}} — \emph{ôbhimān} — \textquotedblleft{}hurt pride; a loving sulk\textquotedblright{}.

\begin{quote}
\bn{বোনের} \bn{জন্মদিনে} \bn{আমি} \bn{ফোন} \bn{করলাম} \bn{না}, \bn{তাই} \bn{তার} \bn{অভিমান} \bn{হল।}
\end{quote}

\begin{quote}
\emph{boner jônmôdine āmi phon kôrlām nā, tāi tār ôbhimān hôlô.}
\end{quote}

\begin{quote}
I didn't phone on my sister's birthday, so she was hurt.
\end{quote}

\bn{অভিমান} is felt between people who love each other: the sulk of a child, a sister or a wife, waiting to be coaxed round. \bn{অভিমানী}, quick to feel it.

\begin{cousinweb}
\textbf{\bn{অভিমান}} (\emph{ôbhimān}) is a tatsama, Sanskrit \emph{abhimāna}, \textquotedblleft{}a high opinion of oneself, pride\textquotedblright{} (Monier-Williams); Bengali turned it into the wounded pride of someone slighted by a person they love.
\end{cousinweb}

\begin{grammarlens}[title={What you've built}]
One more word for describing what you lived and felt.
\end{grammarlens}

\subsection*{Your turn}
\begin{itemize}
\raggedright
  \item \emph{Say it:} \emph{ôbhimān}
  \item \emph{Say it:} \emph{ôbhimān}, once more
  \item \emph{Say it:} \emph{boner jônmôdine āmi phon kôrlām nā, tāi tār ôbhimān hôlô.} — I didn't phone on my sister's birthday, so she was hurt.
  \item \emph{Recall:} say the Bengali for to seethe, then the Bengali for to grumble, then say \emph{ôbhimān} again
\end{itemize}

\begin{tcolorbox}[breakable,colback=teal!4,colframe=teal!35!black,title={Before you move on}]
What does \bn{অভিমান} mean? (\textquotedblleft{}Hurt pride; a loving sulk\textquotedblright{}.) Say it once more.
\end{tcolorbox}
\section[\texorpdfstring{\bn{অনুযোগ} (ônujog)}{অনুযোগ (ônujog)}]{\bn{অনুযোগ} (ônujog) — a gentle complaint; a reproach (to someone close)}

\label{lesson:BN-C444-onujog}



\begin{quote}
Before the new one: say the Bengali for to grumble, then the Bengali for hurt pride.
\end{quote}

\subsection*{You'll want to know: \bn{অনুযোগ}}
\textbf{\bn{অনুযোগ}} — \emph{ônujog} — \textquotedblleft{}a gentle complaint; a reproach (to someone close)\textquotedblright{}.

\begin{quote}
\bn{মীরার} \bn{অনুযোগ}, \bn{অরুণ} \bn{সবসময়} \bn{দেরি} \bn{করে} \bn{আসে।}
\end{quote}

\begin{quote}
\emph{mīrār ônujog, ôrun sôbsômôy deri kôre āse.}
\end{quote}

\begin{quote}
Mira's complaint is that Arun always turns up late.
\end{quote}

\bn{অনুযোগ} \bn{করা}, to complain to someone close. \bn{অভিযোগ} is a formal complaint or a charge; \bn{অনুযোগ} stays inside the family and among friends.

\begin{cousinweb}
\textbf{\bn{অনুযোগ}} (\emph{ônujog}) is a tatsama, Sanskrit \emph{anuyoga}, which Monier-Williams gives for questioning and also for reproof; Bengali uses it for the milder complaint made to someone you are close to.
\end{cousinweb}

\begin{grammarlens}[title={What you've built}]
One more word for describing what you lived and felt.
\end{grammarlens}

\subsection*{Your turn}
\begin{itemize}
\raggedright
  \item \emph{Say it:} \emph{ônujog}
  \item \emph{Say it:} \emph{ônujog}, once more
  \item \emph{Say it:} \emph{mīrār ônujog, ôrun sôbsômôy deri kôre āse.} — Mira's complaint is that Arun always turns up late.
  \item \emph{Recall:} say the Bengali for to grumble, then the Bengali for hurt pride, then say \emph{ônujog} again
\end{itemize}

\begin{tcolorbox}[breakable,colback=teal!4,colframe=teal!35!black,title={Before you move on}]
What does \bn{অনুযোগ} mean? (\textquotedblleft{}A gentle complaint; a reproach (to someone close)\textquotedblright{}.) Say it once more.
\end{tcolorbox}
\section[\texorpdfstring{\bn{মুখ} \bn{ভার} \bn{করা} (mukh bhār kôrā)}{মুখ ভার করা (mukh bhār kôrā)}]{\bn{মুখ} \bn{ভার} \bn{করা} (mukh bhār kôrā) — to sulk; to put on a sullen face}

\label{lesson:BN-C444-mukh-bhar-kora}



\begin{quote}
Before the new one: say the Bengali for hurt pride, then the Bengali for a gentle complaint.
\end{quote}

\subsection*{You'll want to know: \bn{মুখ} \bn{ভার} \bn{করা}}
\textbf{\bn{মুখ} \bn{ভার} \bn{করা}} — \emph{mukh bhār kôrā} — \textquotedblleft{}to sulk; to put on a sullen face\textquotedblright{}.

\begin{quote}
\bn{বকা} \bn{খেয়ে} \bn{ছোট} \bn{বোন} \bn{মুখ} \bn{ভার} \bn{করে} \bn{বসে} \bn{আছে।}
\end{quote}

\begin{quote}
\emph{bôkā kheye chhoṭo bon mukh bhār kôre bôse āchhe.}
\end{quote}

\begin{quote}
After being scolded, my little sister is sitting there sulking.
\end{quote}

Past (he/she …): \bn{তিনি} \bn{মুখ} \bn{ভার} \bn{করলেন}, \emph{tini mukh bhār kôrlen}.

\begin{cousinweb}
\textbf{\bn{মুখ} \bn{ভার} \bn{করা}} (\emph{mukh bhār kôrā}) is \bn{মুখ} (taught), the face, with \emph{bhār}, \textquotedblleft{}weight\textquotedblright{}, from Sanskrit \emph{bhāra}, \textquotedblleft{}burden, weight\textquotedblright{} (CDIAL 9459), and \emph{kôrā}, \textquotedblleft{}to do\textquotedblright{} (CDIAL 2814): to make the face heavy.
\end{cousinweb}

\begin{grammarlens}[title={What you've built}]
One more word for describing what you lived and felt.
\end{grammarlens}

\subsection*{Your turn}
\begin{itemize}
\raggedright
  \item \emph{Say it:} \emph{mukh bhār kôrā}
  \item \emph{Say it:} \emph{mukh bhār kôrā}, once more
  \item \emph{Say it:} \emph{bôkā kheye chhoṭo bon mukh bhār kôre bôse āchhe.} — After being scolded, my little sister is sitting there sulking.
  \item \emph{Recall:} say the Bengali for hurt pride, then the Bengali for a gentle complaint, then say \emph{mukh bhār kôrā} again
\end{itemize}

\begin{tcolorbox}[breakable,colback=teal!4,colframe=teal!35!black,title={Before you move on}]
What does \bn{মুখ} \bn{ভার} \bn{করা} mean? (\textquotedblleft{}To sulk; to put on a sullen face\textquotedblright{}.) Say it once more.
\end{tcolorbox}
\section[\texorpdfstring{\bn{ক্ষমা} \bn{করা} (khômā kôrā)}{ক্ষমা করা (khômā kôrā)}]{\bn{ক্ষমা} \bn{করা} (khômā kôrā) — to forgive}

\label{lesson:BN-C444-khoma-kora}



\begin{quote}
Before the new one: say the Bengali for a gentle complaint, then the Bengali for to sulk.
\end{quote}

\subsection*{You'll want to know: \bn{ক্ষমা} \bn{করা}}
\textbf{\bn{ক্ষমা} \bn{করা}} — \emph{khômā kôrā} — \textquotedblleft{}to forgive\textquotedblright{}.

\begin{quote}
\bn{আমার} \bn{ভুলের} \bn{জন্য} \bn{বন্ধু} \bn{আমাকে} \bn{ক্ষমা} \bn{করল।}
\end{quote}

\begin{quote}
\emph{āmār bhuler jônyô bôndhu āmāke khômā kôrlô.}
\end{quote}

\begin{quote}
My friend forgave me for my mistake.
\end{quote}

Past (he/she …): \bn{তিনি} \bn{ছেলেকে} \bn{ক্ষমা} \bn{করলেন}, \emph{tini chheleke khômā kôrlen}.

\bn{ক্ষমা} \bn{চাওয়া}, to ask forgiveness: \bn{আমি} \bn{ক্ষমা} \bn{চাইলাম}.

\begin{cousinweb}
\textbf{\bn{ক্ষমা} \bn{করা}} (\emph{khômā kôrā}) is Sanskrit \emph{kṣamā}, \textquotedblleft{}patience, forbearance\textquotedblright{} (Monier-Williams), with \emph{kôrā}, \textquotedblleft{}to do\textquotedblright{} (CDIAL 2814).
\end{cousinweb}

\begin{grammarlens}[title={What you've built}]
One more word for describing what you lived and felt.
\end{grammarlens}

\subsection*{Your turn}
\begin{itemize}
\raggedright
  \item \emph{Say it:} \emph{khômā kôrā}
  \item \emph{Say it:} \emph{khômā kôrā}, once more
  \item \emph{Say it:} \emph{āmār bhuler jônyô bôndhu āmāke khômā kôrlô.} — My friend forgave me for my mistake.
  \item \emph{Recall:} say the Bengali for a gentle complaint, then the Bengali for to sulk, then say \emph{khômā kôrā} again
\end{itemize}

\begin{tcolorbox}[breakable,colback=teal!4,colframe=teal!35!black,title={Before you move on}]
What does \bn{ক্ষমা} \bn{করা} mean? (\textquotedblleft{}To forgive\textquotedblright{}.) Say it once more.
\end{tcolorbox}
\section[\texorpdfstring{\bn{মনোমালিন্য} (mônomālinyô)}{মনোমালিন্য (mônomālinyô)}]{\bn{মনোমালিন্য} (mônomālinyô) — ill feeling; bad blood (between people)}

\label{lesson:BN-C444-monomalinyo}



\begin{quote}
Before the new one: say the Bengali for to sulk, then the Bengali for to forgive.
\end{quote}

\subsection*{You'll want to know: \bn{মনোমালিন্য}}
\textbf{\bn{মনোমালিন্য}} — \emph{mônomālinyô} — \textquotedblleft{}ill feeling; bad blood (between people)\textquotedblright{}.

\begin{quote}
\bn{জমি} \bn{নিয়ে} \bn{দুই} \bn{ভাইয়ের} \bn{মনোমালিন্য} \bn{হল।}
\end{quote}

\begin{quote}
\emph{jômi niye dui bhāiyer mônomālinyô hôlô.}
\end{quote}

\begin{quote}
The two brothers fell out over land.
\end{quote}

\begin{cousinweb}
\textbf{\bn{মনোমালিন্য}} (\emph{mônomālinyô}) joins Sanskrit \emph{manas}, \textquotedblleft{}mind\textquotedblright{} (CDIAL 9822), in its compound form \emph{mano-}, and \emph{mālinya}, \textquotedblleft{}dirtiness, darkness\textquotedblright{} (Monier-Williams): a mind clouded against someone.
\end{cousinweb}

\begin{grammarlens}[title={What you've built}]
One more word for describing what you lived and felt.
\end{grammarlens}

\subsection*{Your turn}
\begin{itemize}
\raggedright
  \item \emph{Say it:} \emph{mônomālinyô}
  \item \emph{Say it:} \emph{mônomālinyô}, once more
  \item \emph{Say it:} \emph{jômi niye dui bhāiyer mônomālinyô hôlô.} — The two brothers fell out over land.
  \item \emph{Recall:} say the Bengali for to sulk, then the Bengali for to forgive, then say \emph{mônomālinyô} again
\end{itemize}

\begin{tcolorbox}[breakable,colback=teal!4,colframe=teal!35!black,title={Before you move on}]
What does \bn{মনোমালিন্য} mean? (\textquotedblleft{}Ill feeling; bad blood (between people)\textquotedblright{}.) Say it once more.
\end{tcolorbox}
